\documentclass[sigconf, nonacm]{acmart}
\usepackage{pvldb}

\usepackage{amsmath,amsfonts}
\usepackage{booktabs}
\usepackage{multirow}
\usepackage{array}
\usepackage{microtype}
\usepackage{graphicx}
\usepackage{enumitem}
\usepackage{xcolor}

\newtheorem{theorem}{Theorem}
\newtheorem{lemma}{Lemma}
\newtheorem{proposition}{Proposition}
\newtheorem{corollary}{Corollary}
\newtheorem{definition}{Definition}
\newtheorem{remark}{Remark}

\emergencystretch=3em

\begin{document}

\title[SDG-MH: Multi-Horizon User-Time DP for Perpetual LLM Analytics]{SDG-MH: Certified Calendar-Planned Multi-Horizon User-Time Differential Privacy for Perpetual LLM Analytics}

%% Single-blind review requirement for PVLDB 2027: explicit placeholder block (never fabricating identities)
\author{[Author Name(s) Placeholder -- To Be Completed Prior to Submission]}
\affiliation{%
  \institution{[Institution / Department Affiliation Placeholder]}
  \city{[City Placeholder]}
  \state{[State/Province]}
  \country{[Country Placeholder]}
}
\email{author.placeholder@institution-placeholder.edu}

\begin{abstract}
Enterprise Large Language Model (LLM) telemetry platforms must continuously publish daily operational metrics (total requests, failure counts, provider usage) alongside rolling multi-window distinct-user prevalence (7-day, 30-day, and 90-day active users) over an unbounded operational lifetime. Existing continual differential privacy (DP) mechanisms---such as binary tree aggregation, matrix/factorization mechanisms, and $w$-event DP---either bound privacy under a single window length $w$, assume linear additive prefix queries across a fixed horizon $T$, or force artificial epoch resets. None simultaneously enforce a multi-knot privacy curve $\bar{\rho}(h)$ across short and long user-time horizons ($h \in \{1, 7, 30, 90, 365\}$ days and an affine long-term rate for $h \ge 365$) while supporting non-additive rolling distinct-user counts.

We present \textbf{SDG-MH} (\emph{Sealed Daily Gaussian with Multi-Horizon Calendar Planning}), a perpetual release architecture and certified optimization framework. First, we formalize \textbf{Multi-Horizon User-Time zCDP} via exact window-intersection charging between each release's data-dependency support $B_r$ and any contiguous $h$-day user-time interval $I_{s,h} = [s, s+h-1]$, enforced by an atomic reserve-then-publish-then-commit state machine in our \texttt{UnifiedPrivacyLedger}. Second, we prove a \textbf{210-Day Periodicity Theorem} and an \textbf{Affine-Tail Finite-to-Infinite Horizon Reduction Theorem} showing that verifying $s \in [0, 209]$ and $h \in [1, 575]$ alongside an asymptotic slope inequality $\text{slope}_{\infty} \le \beta_{\text{tail}}$ suffices to guarantee privacy over $t \in [1, \infty)$ without assuming global concavity of $\bar{\rho}(h)$. Third, using the Orbit--Stabilizer Theorem and Chinese Remainder Theorem (CRT) on $\Phi = \mathbb{Z}_7 \times \mathbb{Z}_7 \times \mathbb{Z}_{30}$ under the common time-shift group $\mathbb{Z}_{210}$, we prove that all $1{,}470$ phase combinations (and the $210$ slice with $\phi_7=0$) collapse into \textbf{7 canonical phase equivalence classes} represented by $\{(0, d, 0) : d \in \mathbb{Z}_7\}$, reducing joint phase-and-budget planning to 7 convex subproblems solved in $0.08\text{--}0.19\,\text{s}$ and certified via 100-digit directed-rounding (\texttt{DECIMAL100}) Lagrangian duality. Across a 420-day evaluation (10 independent data seeds $\times$ 10 paired CRN Monte Carlo trials), we report paired statistical comparisons, staleness-versus-accuracy trade-offs, and post-processing-consistent comparisons against binary-tree and $w$-event baselines.
\end{abstract}

\maketitle
\vldbtopmatter

\section{Introduction}
\label{sec:intro}

Organizations deploying internal and customer-facing LLM assistants rely on continuous telemetry dashboards to monitor daily request volume, failure rates, provider routing shares, and rolling active-user adoption (7-day, 30-day, and 90-day unique active users). Because individual prompts and interaction patterns reflect sensitive employee workflows, telemetry releases must satisfy rigorous differential privacy (DP) guarantees~\cite{dwork2006calibrating}. However, operating an enterprise analytics pipeline indefinitely introduces four fundamental challenges that are not jointly resolved by existing continual DP frameworks:

\begin{enumerate}[leftmargin=*, nosep, topsep=2pt]
    \item \textbf{Perpetual Operation Without Artificial Epoch Resets:} Enterprise telemetry has no terminal horizon $T$. Standard epoch-based systems partition time into fixed blocks (e.g., 90-day epochs), which either abruptly change noise scales at epoch boundaries, fail to bound windows straddling two epochs, or exhaust a static lifetime budget.
    \item \textbf{Multi-Horizon Privacy Commitments vs.\ Single-$w$ Bounds:} A single daily cap ($h=1$) allows rapid linear privacy degradation over a year, whereas a single annual cap ($h=365$) permits bursty mechanisms to spend an entire month's budget on a single day. Regulators and privacy engineers require simultaneous bounds $\bar{\rho}(h)$ across multiple user-time horizons ($h \in \{1, 7, 30, 90, 365\}$ days) plus a bounded asymptotic rate $\beta_{\text{tail}}$ for $h > 365$.
    \item \textbf{Mixed Additive and Non-Additive Rolling Distinct Queries:} While daily request and failure counts are additive across days, rolling $W$-day distinct-user counts ($W \in \{7, 30, 90\}$) are non-additive across days because a user active on multiple days within $[t-W+1, t]$ contributes at most $1$ to the window distinct count per cell. Binary trees and Toeplitz/banded matrix factorizations designed for linear prefix sums cannot reconstruct rolling distinct users from daily noisy leaves.
    \item \textbf{Calendar Collision Peaks and Floating-Point Certification Hazards:} Periodic rolling distinct releases collide on specific calendar days. Naive phase alignment $\boldsymbol{\phi} = (0,0,0)$ creates sharp multi-horizon privacy spikes that render historical allocations infeasible ($L_{\max}(2) > \bar{\rho}(2)$). Furthermore, standard floating-point nonlinear solvers (e.g., \texttt{SLSQP}) routinely terminate with $\sim 10^{-12}$ constraint violations, and naive float64 dual deductions $(1 - 10^{-14})$ fail under catastrophic cancellation.
\end{enumerate}

\noindent\textbf{Contributions and Explicit Novelty Scope.}
To prevent overclaiming, Table~\ref{tab:novelty_matrix} classifies every component of SDG-MH into four strict categories: (I) new privacy problem formulation, (II) new mathematical theorems and reductions, (III) direct applications of known identities, and (IV) system and numerical engineering integration. Concretely:
\begin{itemize}[leftmargin=*, nosep, topsep=2pt]
    \item \textbf{Multi-Horizon User-Time zCDP \& Window-Intersection Accounting (Section~\ref{sec:model}):} We formulate perpetual multi-horizon user-time zCDP by charging every release $r$ with data-dependency window $B_r \subseteq \mathbb{Z}_{\ge 1}$ whenever $B_r \cap [s, s+h-1] \neq \emptyset$. We prove via explicit counterexample that release-timestamp-only charging undercounts multi-day distinct windows by up to $2\times$, and enforce atomic admission via a reserve-then-publish-then-commit protocol with a unified $\varepsilon$-drift bound across all commitment curves and horizons (Lemma~\ref{lem:drift}).
    \item \textbf{210-Day Periodicity \& 575-Day Affine-Tail Finite Reduction (Section~\ref{sec:theory}):} For periodic schedules with step sizes $s_f \in \{7, 7, 30\}$ ($P = \operatorname{lcm}(7,7,30) = 210$) and commitment curves that are affine for $h \ge H_0 = 365$, we prove that infinite-horizon security over $(s, h) \in \mathbb{Z}_{\ge 0} \times \mathbb{Z}_{\ge 1}$ reduces exactly to verifying $s \in [0, 209]$ and $h \in [1, H_0 + P] = [1, 575]$ plus $\text{slope}_{\infty} \le \beta_{\text{tail}}$---without assuming global concavity of $\bar{\rho}(h)$ (which fails at $h=30$ for both commitments).
    \item \textbf{CRT 7-Orbit Phase Equivalence \& Certified \texttt{DECIMAL100} Joint Planning (Section~\ref{sec:planner}):} Using the Orbit--Stabilizer Theorem and Chinese Remainder Theorem on $\Phi = \mathbb{Z}_7 \times \mathbb{Z}_7 \times \mathbb{Z}_{30}$ under the common shift group $\mathbb{Z}_{210}$, we prove that all $1{,}470$ phase tuples (and the $210$ tuples with $\phi_7=0$) partition into \textbf{7 canonical equivalence classes} with representatives $\{(0, d, 0) : d \in \mathbb{Z}_7\}$ (Proposition~\ref{prop:orbits}), cutting unpinned joint planning from 210 redundant solves ($1.28\text{--}5.79\,\text{s}$) to 7 canonical solves ($0.08\text{--}0.19\,\text{s}$) with identical certified primal and 100-digit directed-rounding (\texttt{DECIMAL100}) dual bounds. We also formalize the three operational conditions that break $\mathbb{Z}_{210}$ shift equivalence.
    \item \textbf{420-Day Cross-Period Paired Evaluation (Section~\ref{sec:experiments}):} Across 10 independent data seeds $\times$ 10 paired CRN trials over $T=420=2P$ days, we report seed-level paired confidence intervals, Holm--Bonferroni corrections, same-query-date accuracy vs.\ staleness decomposition, and why abstract variance reduction ($-20.05\%$) can coexist with higher Distinct MAPE ($+4.03\%$ to $+4.45\%$) due to negative truncation bias.
\end{itemize}

\begin{table*}[t]
\centering
\caption{Strict four-way classification of SDG-MH claims against prior continual DP and privacy accounting literature.}
\label{tab:novelty_matrix}
\resizebox{\textwidth}{!}{%
\begin{tabular}{@{}p{3.6cm}p{2.9cm}p{5.4cm}p{5.2cm}@{}}
\toprule
\textbf{Paper Component} & \textbf{Classification} & \textbf{What Is New in SDG-MH} & \textbf{What Must NOT Be Claimed as a Breakthrough} \\ \midrule
Multi-Horizon User-Time zCDP ($\bar{\rho}(h)$ at $h \in \{1,7,30,90,365\}$ + affine tail) & \textbf{(I) New Problem Formulation} & Generalizes single-window $w$-event DP~\cite{kellaris2014wevent} to simultaneous multi-knot user-time bounds over $t \in [1, \infty)$ with exact window-overlap charging $B_r \cap I_{s,h} \neq \emptyset$. & Does not invent a new divergence family beyond standard zCDP~\cite{bun2016concentrated}; does not replace lifetime user-level DP if full-history immunity ($h=\infty, \beta_{\text{tail}}=0$) is required. \\ \addlinespace[2pt]
210-Day Periodicity (Thm.~\ref{thm:periodicity}) \& 575-Day Affine-Tail Reduction (Thm.~\ref{thm:reduction}) & \textbf{(II) New Mathematical Theorems} & Exact finite reduction of $\forall (s,h) \in \mathbb{Z}_{\ge 0} \times \mathbb{Z}_{\ge 1}$ to $s \in [0, 209], h \in [1, 575]$ + asymptotic slope check without requiring global concavity of $\bar{\rho}(h)$. & For a fixed single window $w$, periodic intersection counting is elementary modular arithmetic; the novelty is the exact non-concave multi-knot + affine-tail reduction. \\ \addlinespace[2pt]
CRT 7-Orbit Phase Reduction (Prop.~\ref{prop:orbits}) \& Symmetry-Breaking Conditions & \textbf{(II) New Mathematical Proposition} & Proves via Orbit--Stabilizer and CRT ($\delta = (120\phi_7 + 91\phi_{90}) \bmod 210$) that $|\Phi/\mathbb{Z}_{210}| = 1470/210 = 7$ (exposing a $30\times$ redundancy in naive $(0,\phi_{30},\phi_{90})$ enumeration). & Elementary finite group action / CRT once formulated; must not be claimed as an NP-hard combinatorial optimization breakthrough. \\ \addlinespace[2pt]
Daily Independent Gaussian + Window-Level Distinct Gaussian & \textbf{(III) Direct Application of Known Mechanisms} & Combines daily additive sealing with periodic window-level distinct-user Gaussian releases under unified multi-horizon zCDP accounting. & Does \textbf{not} invent a new DP noise mechanism; for pure additive prefix queries, binary trees~\cite{dwork2010continual,chan2011private} and Toeplitz/banded MF~\cite{choquette2023banded,henzinger2024group,mcmahan2025blt} achieve $O(\log^3 h)$ vs.\ SDG-MH's $O(h)$ variance. \\ \addlinespace[2pt]
KKT Closed-Form Single-Constraint \& \texttt{DECIMAL100} Dual Certificate (Prop.~\ref{prop:dual}) & \textbf{(III) Known Convex Duality + (IV) Certified Engineering} & Eliminates float64 catastrophic cancellation in $g^\top \hat{\boldsymbol{v}} - \boldsymbol{\lambda}^\top \boldsymbol{b}$ via 100-digit directed-rounding enclosure arithmetic (\texttt{ROUND\_FLOOR}/\texttt{ROUND\_CEILING}). & Separable convex programming $\min \sum c_j/v_j$ s.t.\ $A\boldsymbol{v} \le \boldsymbol{b}$ is classical; numerical SLSQP + dual enclosure is \texttt{NEAR\_GLOBAL\_OPTIMAL\_CERTIFIED}, not symbolic proof. \\ \addlinespace[2pt]
Atomic Two-Phase Ledger \&Conservative Budget Shrinkage & \textbf{(IV) System \& Numerical Engineering} & Prevents uncharged partial-write leaks (\texttt{RESERVED} $\to$ persist $\to$ \texttt{COMMITTED}) and repairs $-2.99 \times 10^{-12}$ SLSQP violations to strictly positive slack. & Engineering safety pattern (write-ahead reservation + inward ray scaling), not a theoretical cryptography primitive. \\ \bottomrule
\end{tabular}%
}
\end{table*}

\section{Related Work and Prior-Art Positioning}
\label{sec:related}

\noindent\textbf{Continual Observation and Binary Tree Mechanisms.}
Dwork et al.~\cite{dwork2010continual} and Chan et al.~\cite{chan2011private} introduced the binary tree mechanism for continual prefix sums over horizon $T$, achieving $O(\text{polylog } T)$ error by decomposing intervals into $O(\log T)$ dyadic nodes. Recent work studies growing databases and user-level contribution bounding~\cite{cummings2018growing,amin2019bounding} as well as pan-private/streaming frequency and distinct estimation~\cite{mir2011panprivate,epasto2023frequency,ghazi2023distinct,jain2023turnstile}. However, binary trees require either a predetermined tree height $L = \lceil \log_2 T \rceil$ or periodic tree restarts, and cannot compute rolling $W$-day distinct-user counts from daily additive leaves because set union is non-linear.

\noindent\textbf{Matrix Mechanisms and Banded/Toeplitz Factorizations.}
The matrix mechanism~\cite{li2015matrix} factorizes a workload matrix $W = B C$, adding noise to $C x$ and reconstructing via $B$. Denisov et al.~\cite{denisov2022improved}, Fichtenberger et al.~\cite{fichtenberger2023constant}, Choquette-Choo et al.~\cite{choquette2023multiepoch,choquette2023banded}, Henzinger and Upadhyay~\cite{henzinger2024group}, and McMahan and Pillutla~\cite{mcmahan2025blt} developed optimal Toeplitz, buffered linear Toeplitz (BLT), and $b$-banded matrix factorizations for continual linear queries. While banded factorizations enjoy simple sliding-window sensitivity under $w$-event DP, they apply to \emph{linear additive} queries over daily vectors and do not solve periodic multi-window distinct-user scheduling under a multi-knot user-time commitment curve $\bar{\rho}(h)$.

\noindent\textbf{$w$-Event DP vs.\ Multi-Horizon User-Time DP.}
Kellaris et al.~\cite{kellaris2014wevent} proposed $w$-event DP to bridge event-level and user-level privacy over any $w$ consecutive steps, alongside decayed-sum stream models~\cite{bolot2013decayed}. Standard $w$-event DP enforces a single scalar budget $\varepsilon$ for a single fixed window size $w$. By contrast, SDG-MH enforces a multi-horizon commitment curve $\{\bar{\rho}(h)\}_{h \ge 1}$ across $h \in \{1, 7, 30, 90, 365\}$ and all $h > 365$ simultaneously, charging each release $r$ whenever its data-dependency window $B_r$ intersects the protected user-time interval $I_{s,h} = [s, s+h-1]$.

\noindent\textbf{Meaning of ``No Epochs, No Expiry, Perpetual Analytics''.}
In SDG-MH, \emph{No Epochs} means the system eliminates hard 90-day epoch boundaries that reset or block queries; \emph{No Expiry} means historical sealed daily/window outputs remain permanently queryable without invalidation; and \emph{Perpetual Analytics} means the schedule operates over $t \in [1, \infty)$ while satisfying $\text{Exposure}(s, h) \le \bar{\rho}(h)$ for every start day $s \ge 0$ and every finite window length $h \ge 1$. Under a strictly positive tail rate $\beta_{\text{tail}} > 0$, cumulative exposure grows linearly as $\Theta(h)$ for a hypothetical user active forever ($h \to \infty$); bounded $h$-day protection holds for every finite $h$.

\section{System Model and Multi-Horizon User-Time zCDP}
\label{sec:model}

\subsection{Telemetry Schema, Sensitivity, and Truncation}
Time is partitioned into discrete UTC days $t \in \mathbb{Z}_{\ge 1}$. Each telemetry event carries fields \texttt{interaction\_id}, \texttt{user\_id}, \texttt{timestamp}, \texttt{provider}, \texttt{task}, and \texttt{success}, mapped onto $M = 3 \times 3 = 9$ provider-by-task cells. SDG-MH publishes two classes of statistics:
\begin{enumerate}[leftmargin=*, nosep, topsep=2pt]
    \item \textbf{Daily Additive Releases ($t \in \mathbb{Z}_{\ge 1}$):} For each day $t$, each user $u$'s daily cell request vector $r_{u,t} \in \mathbb{R}_{\ge 0}^9$ and failure vector $f_{u,t} \in \mathbb{R}_{\ge 0}^9$ ($0 \le f_{u,t,m} \le r_{u,t,m}$) are clipped to $\ell_2$-balls of radius $C_{\text{req}} = 50.0$ and $C_{\text{fail}} = 10.0$, respectively, with data-dependency window $B_t^{\text{daily}} = \{t\}$. Under isotropic Gaussian noise $\mathcal{N}(0, \sigma_{\text{day}}^2 I_{18})$ on the concatenated vector $[r_t; f_t]$, the daily zCDP cost is $d = \rho_{\text{req}} + \rho_{\text{fail}} = \frac{C_{\text{req}}^2 + C_{\text{fail}}^2}{2 \sigma_{\text{day}}^2}$ with optimal sensitivity-squared split $\eta = \frac{C_{\text{req}}^2}{C_{\text{req}}^2 + C_{\text{fail}}^2} = \frac{2500}{2600} = \frac{25}{26}$.
    \item \textbf{Periodic Rolling Distinct-User Releases:} For each window family $f \in \{7, 30, 90\}$ with window length $W_f \in \{7, 30, 90\}$, step period $s_f \in \{7, 7, 30\}$, and phase $\phi_f \in \mathbb{Z}_{s_f}$, a release is triggered on every day $t \ge 1$ satisfying $t \equiv \phi_f \pmod{s_f}$ with data-dependency window $B_{f,t} = [\max(1, t - W_f + 1), t]$. Within $B_{f,t}$, each user $u$ has a binary cell-presence vector $b_{u,B_{f,t}} \in \{0,1\}^9$ plus 2 overall/margin indicators clipped to $\ell_2$-sensitivity $\Delta_2^{\text{dist}} = \sqrt{11}$, consuming zCDP budget $x_f = \frac{11}{2 \sigma_f^2}$ per release.
\end{enumerate}

\subsection{Window-Intersection Accounting vs.\ Timestamp Charging}
\begin{definition}[Multi-Horizon User-Time Adjacency and Exposure]
\label{def:exposure}
For any start day $s \in \mathbb{Z}_{\ge 0}$ and horizon length $h \in \mathbb{Z}_{\ge 1}$, let $I_{s,h} = [s, s+h-1]$ denote a contiguous $h$-day user-time window. Two telemetry streams $D, D'$ are $(s, h)$-user-time adjacent ($D \sim_{s,h} D'$) if they differ only in the interactions of a single user $u^\star$ occurring on days $t \in I_{s,h}$. Given a collection of releases $\mathcal{R}$ where release $r \in \mathcal{R}$ accesses data only from days $B_r \subseteq \mathbb{Z}_{\ge 1}$ and is $\rho_r$-zCDP with respect to adding/removing one user's contributions on $B_r$, the $(s, h)$-exposure is:
\begin{equation}
\label{eq:exposure}
\mathrm{Exposure}(s, h; \mathcal{R}) \;=\; \sum_{r \in \mathcal{R} \,:\, B_r \cap I_{s,h} \neq \emptyset} \rho_r.
\end{equation}
A schedule satisfies Multi-Horizon User-Time zCDP under commitment curve $\bar{\rho}: \mathbb{Z}_{\ge 1} \to \mathbb{R}_{>0}$ if $\sup_{s \in \mathbb{Z}_{\ge 0}} \mathrm{Exposure}(s, h; \mathcal{R}) \le \bar{\rho}(h)$ for all $h \in \mathbb{Z}_{\ge 1}$.
\end{definition}

\noindent\textbf{Counterexample Against Release-Timestamp Charging.}
Suppose a naive accountant charges a periodic window release $B_{f,t} = [t - W_f + 1, t]$ only when its publication timestamp $t \in I_{s,h}$, i.e., $\tilde{n}_f(s, h) = |\{t \in I_{s,h} : t \equiv \phi_f \pmod{s_f}\}|$. Consider $W_{30} = 30$, $s_{30} = 7$, $\phi_{30} = 0$, and a user active only on day $s=36$ ($h=1$, $I_{36,1} = \{36\}$). No release occurs on day 36 since $36 \not\equiv 0 \pmod 7$, so $\tilde{n}_{30}(36, 1) = 0$. Yet the user's day-36 data directly alters the four releases published on days $t \in \{42, 49, 56, 63\}$ because $36 \in [t-29, t]$ for each such $t$! Similarly, for $h=30$ and $s=36$ ($I_{36,30} = [36, 65]$), timestamp charging counts only $\{42, 49, 56, 63\}$ ($\tilde{n}_{30} = 4$), whereas $I_{36,30}$ actually intersects the 8 release windows at $t \in \{42, 49, 56, 63, 70, 77, 84, 91\}$ ($n_{30}(36, 30) = 8$, a $2\times$ undercount).

\begin{lemma}[Exact Closed-Form Window-Intersection Count]
\label{lem:intersection}
For a periodic family $f$ with window length $W_f \ge 1$, step period $s_f \ge 1$, and phase $\phi_f \in \mathbb{Z}_{s_f}$ on the steady-state timeline, a release at $t = k s_f + \phi_f$ has support $B_f(k) = [k s_f + \phi_f - W_f + 1,\, k s_f + \phi_f]$. For any $I_{s,h} = [s, s+h-1]$, $B_f(k) \cap I_{s,h} \neq \emptyset$ if and only if $k_{\min} \le k \le k_{\max}$, where:
\begin{equation}
\label{eq:k_bounds}
k_{\min} = \left\lceil \frac{s - \phi_f}{s_f} \right\rceil, \qquad k_{\max} = \left\lfloor \frac{s + h + W_f - 2 - \phi_f}{s_f} \right\rfloor,
\end{equation}
yielding $n_f(s, h; \phi_f) = \max(0,\, k_{\max} - k_{\min} + 1)$.
\end{lemma}
\begin{proof}
Two closed integer intervals $[a, b]$ and $[c, d]$ intersect iff $b \ge c$ and $a \le d$. Substituting $a = k s_f + \phi_f - W_f + 1$, $b = k s_f + \phi_f$, $c = s$, and $d = s + h - 1$ gives $k s_f + \phi_f \ge s \iff k \ge (s - \phi_f)/s_f$ and $k s_f + \phi_f - W_f + 1 \le s + h - 1 \iff k \le (s + h + W_f - 2 - \phi_f)/s_f$. Since $k \in \mathbb{Z}$, Eqs.~\eqref{eq:k_bounds} follow immediately.
\end{proof}

\subsection{Atomic Two-Phase Ledger and Unified \texorpdfstring{$\varepsilon$}{epsilon}-Drift Bound}
To prevent privacy leaks under storage exceptions, \texttt{Unified\allowbreak Privacy\allowbreak Ledger} enforces a three-state lifecycle (\texttt{RESERVED} $\to$ \texttt{COMMITTED} or \texttt{ABORTED}) via \texttt{admit\_\allowbreak and\_\allowbreak commit\_\allowbreak batch}: (1) a batch of candidate releases is checked atomically against $\bar{\rho}(h)$ across all impacted windows $I_{s,h}$ and recorded as \texttt{RESERVED} (charging full budget immediately); (2) Gaussian noise is drawn and persisted to disk before updating in-memory cache; and (3) the records transition to \texttt{COMMITTED}. Rollback is strictly restricted to \texttt{RESERVED} records prior to public exposure.

Because runtime ledger admission uses an IEEE-754 float64 tolerance $\tau_{\text{tol}} = 10^{-12}$ (while planner certification uses \texttt{strict\_zero\_tol}=\texttt{True}, i.e., $\tau_{\text{tol}} = 0$), we establish a unified upper bound on the induced $(\varepsilon, \delta)$-DP drift across all commitment curves and horizons.

\begin{lemma}[Unified Conservative $\varepsilon$-Drift Bound Across All Commitments and Horizons]
\label{lem:drift}
Under standard zCDP-to-$(\varepsilon, \delta)$ conversion $\varepsilon(\rho) = \rho + 2\sqrt{\rho \ln(1/\delta)}$ with $\delta \in (0, 1)$, if an admitted schedule satisfies $\mathrm{Exposure}(s, h) \le \rho + \tau_{\text{tol}}$ with $\tau_{\text{tol}} \ge 0$, then for every $\rho > 0$:
\begin{equation}
\label{eq:drift_bound}
\begin{aligned}
\Delta\varepsilon(\rho; \tau_{\text{tol}}) &:= \varepsilon(\rho + \tau_{\text{tol}}) - \varepsilon(\rho) \\
&\le \min\!\left\{\left(1 + \sqrt{\frac{\ln(1/\delta)}{\rho}}\right)\tau_{\text{tol}},\;\; \tau_{\text{tol}} + 2\sqrt{\tau_{\text{tol}}\ln(1/\delta)}\right\}.
\end{aligned}
\end{equation}
For $\delta = 10^{-6}$ ($\ln(1/\delta) = 13.815511$) and $\tau_{\text{tol}} = 10^{-12}$:
\begin{enumerate}[leftmargin=*, nosep, topsep=1pt]
    \item Across all horizons $h \ge 1$ of both source-code commitment curves in Table~\ref{tab:commitments}, since $\bar{\rho}(h)$ is strictly increasing in $h$, the global minimum commitment is achieved at $h=1$ on \texttt{Original\_Planner} with $\bar{\rho}_{\min} = \bar{\rho}_{\mathrm{Orig}}(1) = 0.01419861$ ($\varepsilon(1)=0.90$), yielding $\Delta\varepsilon \le \mathbf{3.2193 \times 10^{-11}}$; on \texttt{Week1\_v2}, $\bar{\rho}_{\min} = \bar{\rho}_{\mathrm{W1v2}}(1) = 0.01746890$ ($\varepsilon(1)=1.00$), yielding $\Delta\varepsilon \le \mathbf{2.9122 \times 10^{-11}}$. (Even at a legacy draft knot $\varepsilon(1)=0.70 \implies \bar{\rho}(1)=0.00894312$, $\Delta\varepsilon \le 4.0273 \times 10^{-11}$.)
    \item At the smallest individual sub-channel allocation ($\rho_{1,\mathrm{add}} = 0.00030204$), $\Delta\varepsilon \le 2.1471 \times 10^{-10}$; at the solver box floor $l_{\min} = 10^{-6}$, $\Delta\varepsilon \le 3.7184 \times 10^{-9}$; and for arbitrary $\rho \ge 0$, the universal non-singular root bound gives $\Delta\varepsilon \le 7.4338 \times 10^{-6}$.
\end{enumerate}
\end{lemma}
\begin{proof}
For $\rho > 0$, $\varepsilon''(\rho) = -\frac{1}{2}\sqrt{\ln(1/\delta)}\rho^{-3/2} < 0$, so $\varepsilon(\rho)$ is strictly concave and bounded above by its first-order Taylor expansion at $\rho$: $\varepsilon(\rho + \tau_{\text{tol}}) - \varepsilon(\rho) \le \varepsilon'(\rho)\tau_{\text{tol}} = (1 + \sqrt{\ln(1/\delta)/\rho})\tau_{\text{tol}}$. Simultaneously, subadditivity $\sqrt{\rho + \tau_{\text{tol}}} - \sqrt{\rho} \le \sqrt{\tau_{\text{tol}}}$ for all $\rho \ge 0$ yields the non-singular bound $\tau_{\text{tol}} + 2\sqrt{\tau_{\text{tol}}\ln(1/\delta)}$. Taking the minimum proves Eq.~\eqref{eq:drift_bound}. Because both commitment curves have positive secant slopes on $[1, 365]$ and $\beta_{\text{tail}} > 0$ on $[365, \infty)$, $\min_{h \ge 1} \bar{\rho}(h) = \bar{\rho}(1)$, giving the exact constants above.
\end{proof}

\begin{table*}[t]
\centering
\caption{Source-code privacy commitment curves (\texttt{calendar\_accounting.py}, $\delta = 10^{-6}$), exact zCDP knot values $\bar{\rho}(h)$, adjacent secant slopes $\Delta\bar{\rho}/\Delta h$, affine tail parameters ($\bar{\rho}(h) = a_{\text{tail}} + \beta_{\text{tail}} h$ for $h \ge 365$), and Lemma~\ref{lem:drift} $\varepsilon$-drift bounds ($\tau_{\text{tol}}=10^{-12}$). Note: all entries are derived directly from \texttt{ORIGINAL\_PLANNER\_CURVE} and \texttt{WEEK1\_V2\_CURVE} in code, resolving a prior draft table discrepancy where legacy $\varepsilon=(0.70, 1.08, 1.49, 2.14, 4.29)$ figures were mistakenly listed alongside \texttt{WEEK1\_V2\_CURVE} slopes.}
\label{tab:commitments}
\resizebox{\textwidth}{!}{%
\begin{tabular}{@{}llcccccccccc@{}}
\toprule
\textbf{Commitment Curve} & \textbf{Quantity} & $h=1$ & $h=7$ & $h=30$ & $h=90$ & $h=365$ & \textbf{Affine Tail} $\beta_{\text{tail}}$ & \textbf{Intercept} $a_{\text{tail}}$ & $\bar{\rho}_{\min}$ ($h \ge 1$) & \textbf{Max $\varepsilon$-Drift} & \textbf{Global Concave?} \\ \midrule
\multirow{3}{*}{\texttt{Original\_Planner}} & Target $\varepsilon(h)$ & $0.9000$ & $1.0800$ & $1.4900$ & $2.2400$ & $4.3900$ & \multirow{3}{*}{$0.00079364$} & \multirow{3}{*}{$0.01267882$} & \multirow{3}{*}{$0.01419861$} & \multirow{3}{*}{$3.2193 \times 10^{-11}$} & \multirow{3}{*}{\textbf{No} ($s_{30\to 90} < s_{90\to 365}$)} \\
 & Exact $\bar{\rho}(h)$ & $0.014199$ & $0.020320$ & $0.038143$ & $0.084106$ & $0.302356$ & & & & & \\
 & Secant $\Delta\bar{\rho}/\Delta h$ & --- & $0.0010202$ & $0.0007749$ & $0.0007660$ & $0.0007936$ & & & & & \\ \midrule
\multirow{3}{*}{\texttt{Week1\_v2}} & Target $\varepsilon(h)$ & $1.0000$ & $1.2000$ & $1.5000$ & $2.2500$ & $4.4000$ & \multirow{3}{*}{$0.00079569$} & \multirow{3}{*}{$0.01321923$} & \multirow{3}{*}{$0.01746890$} & \multirow{3}{*}{$2.9122 \times 10^{-11}$} & \multirow{3}{*}{\textbf{No} ($s_{7\to 30} < s_{30\to 90} < s_{90\to 365}$)} \\
 & Exact $\bar{\rho}(h)$ & $0.017469$ & $0.024984$ & $0.038644$ & $0.084831$ & $0.303646$ & & & & & \\
 & Secant $\Delta\bar{\rho}/\Delta h$ & --- & $0.0012525$ & $0.0005939$ & $0.0007698$ & $0.0007957$ & & & & & \\ \bottomrule
\end{tabular}%
}
\end{table*}

\section{Periodicity and Non-Concave Infinite-Horizon Reduction}
\label{sec:theory}

A central challenge in perpetual privacy accounting is verifying $\mathrm{Exposure}(s, h) \le \bar{\rho}(h)$ across infinitely many pairs $(s, h) \in \mathbb{Z}_{\ge 0} \times \mathbb{Z}_{\ge 1}$. As shown in Table~\ref{tab:commitments}, \emph{neither} \texttt{Original\_Planner} nor \texttt{Week1\_v2} is globally concave in $h$: on \texttt{Original\_Planner}, the secant slope increases from $s_{30\to 90} = 0.0007660$ to $s_{90\to 365} = 0.0007936$; on \texttt{Week1\_v2}, $s_{7\to 30} = 0.0005939 < s_{30\to 90} = 0.0007698 < s_{90\to 365} = 0.0007957$. Thus, any finite-reduction proof that assumes global concavity of $\bar{\rho}(h)$ is invalid. Instead, we exploit only the exact \textbf{affine tail property} of $\bar{\rho}(h)$ for $h \ge H_0 = 365$: $\bar{\rho}(h) = a_{\text{tail}} + \beta_{\text{tail}} h$.

\begin{theorem}[210-Day Joint Periodicity]
\label{thm:periodicity}
Let $P = \operatorname{lcm}(s_7, s_{30}, s_{90}) = \operatorname{lcm}(7, 7, 30) = 210$. For any phase vector $\boldsymbol{\phi} \in \mathbb{Z}_7 \times \mathbb{Z}_7 \times \mathbb{Z}_{30}$ and budget vector $\boldsymbol{v} = (d, x_7, x_{30}, x_{90})^\top \in \mathbb{R}_{\ge 0}^4$, the steady-state exposure $L(s, h; \boldsymbol{\phi}, \boldsymbol{v}) = h d + \sum_{f \in \{7,30,90\}} n_f(s, h; \phi_f) x_f$ satisfies:
\begin{enumerate}[leftmargin=*, nosep, topsep=1pt]
    \item $L(s + P, h; \boldsymbol{\phi}, \boldsymbol{v}) = L(s, h; \boldsymbol{\phi}, \boldsymbol{v})$ for all $s \in \mathbb{Z}_{\ge 0}$ and $h \in \mathbb{Z}_{\ge 1}$; hence $L_{\max}(h; \boldsymbol{\phi}, \boldsymbol{v}) := \sup_{s \ge 0} L(s, h; \boldsymbol{\phi}, \boldsymbol{v})$ equals $\max_{s \in \{0, \dots, 209\}} L(s, h; \boldsymbol{\phi}, \boldsymbol{v})$.
    \item For every $s \in \mathbb{Z}_{\ge 0}$ and $h \in \mathbb{Z}_{\ge 1}$, $n_f(s, h + P; \phi_f) = n_f(s, h; \phi_f) + \frac{P}{s_f}$, and therefore:
    \begin{equation}
    \label{eq:L_shift}
    \begin{aligned}
    L_{\max}(h + P; \boldsymbol{\phi}, \boldsymbol{v}) &= L_{\max}(h; \boldsymbol{\phi}, \boldsymbol{v}) + P \cdot \mathrm{slope}_{\infty}(\boldsymbol{v}), \\
    \mathrm{slope}_{\infty}(\boldsymbol{v}) &:= d + \frac{x_7}{7} + \frac{x_{30}}{7} + \frac{x_{90}}{30}.
    \end{aligned}
    \end{equation}
\end{enumerate}
\end{theorem}
\begin{proof}
Since $s_f \mid P$, shifting $s \mapsto s + P$ shifts $k_{\min}$ and $k_{\max}$ in Eq.~\eqref{eq:k_bounds} by $+P/s_f \in \mathbb{Z}$, leaving $k_{\max} - k_{\min} + 1$ invariant. Similarly, replacing $h \mapsto h + P$ leaves $k_{\min}$ unchanged and increments $k_{\max}$ by $+P/s_f \in \mathbb{Z}$, increasing $L(s, h; \boldsymbol{\phi}, \boldsymbol{v})$ by $P d + \sum_f (P/s_f) x_f = P \cdot \mathrm{slope}_{\infty}(\boldsymbol{v})$ for every $s$. Taking $\max_{s \in [0, 209]}$ yields Eq.~\eqref{eq:L_shift}. At startup ($t \ge 1$), the active release set is a subset of the bi-infinite steady-state lattice, so $\mathrm{Exposure}_{\text{startup}}(s, h) \le L_{\max}(h; \boldsymbol{\phi}, \boldsymbol{v})$.
\end{proof}

\begin{theorem}[575-Day Finite-to-Infinite Horizon Reduction via Affine Tail]
\label{thm:reduction}
Let $\bar{\rho}: \mathbb{Z}_{\ge 1} \to \mathbb{R}_{>0}$ be any commitment curve (not necessarily concave on $[1, H_0]$) that is affine for all $h \ge H_0 = 365$: $\bar{\rho}(h) = a_{\mathrm{tail}} + \beta_{\mathrm{tail}} h$. Let $P = 210$ and $H_{\mathrm{eval}} = H_0 + P = 575$. Then $L_{\max}(h; \boldsymbol{\phi}, \boldsymbol{v}) \le \bar{\rho}(h)$ holds for \textbf{all} $h \in \mathbb{Z}_{\ge 1}$ if and only if:
\begin{equation}
\label{eq:finite_reduction}
\begin{aligned}
\text{(i)}\;\; &\max_{1 \le h \le 575} \bigl(L_{\max}(h; \boldsymbol{\phi}, \boldsymbol{v}) - \bar{\rho}(h)\bigr) \;\le\; 0, \\
\text{(ii)}\;\; &\mathrm{slope}_{\infty}(\boldsymbol{v}) \;=\; d + \frac{x_7}{7} + \frac{x_{30}}{7} + \frac{x_{90}}{30} \;\le\; \beta_{\mathrm{tail}}.
\end{aligned}
\end{equation}
\end{theorem}
\begin{proof}
($\Rightarrow$) Necessity of (i) is immediate. If $\mathrm{slope}_{\infty}(\boldsymbol{v}) > \beta_{\mathrm{tail}}$, then for $k \in \mathbb{Z}_{\ge 1}$, $L_{\max}(H_0 + k P) - \bar{\rho}(H_0 + k P) = L_{\max}(H_0) - \bar{\rho}(H_0) + k P (\mathrm{slope}_{\infty}(\boldsymbol{v}) - \beta_{\mathrm{tail}}) \to +\infty$ as $k \to \infty$, violating feasibility.

($\Leftarrow$) Suppose (i) and (ii) hold. Every $h \in [1, 575]$ is directly covered by (i). For any $h > 575$, write $h = h_0 + k P$ where $k = \lfloor (h - H_0)/P \rfloor \ge 1$ and $h_0 = h - k P \in [H_0, H_0 + P - 1] = [365, 574] \subset [1, 575]$. Because both $h_0 \ge 365$ and $h \ge 365$ lie in the affine tail region of $\bar{\rho}$, we have $\bar{\rho}(h) = \bar{\rho}(h_0) + k P \beta_{\mathrm{tail}}$. By Theorem~\ref{thm:periodicity} and conditions (i)--(ii):
\[
\begin{aligned}
L_{\max}(h; \boldsymbol{\phi}, \boldsymbol{v}) &= L_{\max}(h_0; \boldsymbol{\phi}, \boldsymbol{v}) + k P \, \mathrm{slope}_{\infty}(\boldsymbol{v}) \\
&\le \bar{\rho}(h_0) + k P \beta_{\mathrm{tail}} \;=\; \bar{\rho}(h). \qedhere
\end{aligned}
\]
\end{proof}

\section{7-Orbit Phase Reduction and Certified Joint Planning}
\label{sec:planner}

\subsection{Orbit--Stabilizer and CRT Reduction to 7 Canonical Phase Equivalence Classes}
The discrete phase vector is $\boldsymbol{\phi} = (\phi_7, \phi_{30}, \phi_{90}) \in \Phi := \mathbb{Z}_7 \times \mathbb{Z}_7 \times \mathbb{Z}_{30}$, which has $|\Phi| = 7 \times 7 \times 30 = 1{,}470$ tuples. Previously, fixing only $\phi_7 = 0$ by time-origin invariance left $1 \times 7 \times 30 = 210$ candidate tuples $(0, \phi_{30}, \phi_{90})$. We prove that this 210-tuple slice contains an exact $30\times$ redundancy: under the full common time-shift group $G = \mathbb{Z}_{210}$, the entire $1{,}470$-element space $\Phi$ partitions into \textbf{only 7 canonical equivalence classes}.

\begin{proposition}[Orbit--Stabilizer and CRT Reduction to 7 Canonical Phase Orbits]
\label{prop:orbits}
Let the cyclic group $G = \mathbb{Z}_{210}$ act on $\Phi = \mathbb{Z}_7 \times \mathbb{Z}_7 \times \mathbb{Z}_{30}$ by shifting the global calendar origin by $\delta \in \mathbb{Z}_{210}$:
\begin{equation}
\label{eq:group_action}
\begin{aligned}
\delta \cdot (\phi_7, \phi_{30}, \phi_{90}) \;=\; \bigl(&(\phi_7 - \delta) \bmod 7,\;\; (\phi_{30} - \delta) \bmod 7, \\
&(\phi_{90} - \delta) \bmod 30\bigr).
\end{aligned}
\end{equation}
Then:
\begin{enumerate}[leftmargin=*, nosep, topsep=1pt]
    \item \textbf{Free Action and 7 Orbits:} The stabilizer of every $\boldsymbol{\phi} \in \Phi$ is trivial, $\mathrm{Stab}_{\mathbb{Z}_{210}}(\boldsymbol{\phi}) = \{0\}$. By the Orbit--Stabilizer Theorem, every orbit $\mathcal{O}(\boldsymbol{\phi})$ has cardinality $|\mathcal{O}(\boldsymbol{\phi})| = 210$, partitioning $\Phi$ into $|\Phi / \mathbb{Z}_{210}| = 1{,}470 / 210 = \mathbf{7}$ disjoint orbits.
    \item \textbf{Constructive CRT Canonicalization:} Since $\gcd(7, 30) = 1$ with orthogonal idempotents $e_7 = 120 \equiv 1 \pmod 7,\, 120 \equiv 0 \pmod{30}$ and $e_{30} = 91 \equiv 0 \pmod 7,\, 91 \equiv 1 \pmod{30}$, for any $\boldsymbol{\phi} = (\phi_7, \phi_{30}, \phi_{90}) \in \Phi$, the unique shift $\delta(\boldsymbol{\phi}) \in \mathbb{Z}_{210}$ given by
    \begin{equation}
    \label{eq:crt_shift}
    \delta(\boldsymbol{\phi}) \;=\; (120\,\phi_7 + 91\,\phi_{90}) \bmod 210
    \end{equation}
    satisfies $\delta(\boldsymbol{\phi}) \equiv \phi_7 \pmod 7$ and $\delta(\boldsymbol{\phi}) \equiv \phi_{90} \pmod{30}$, mapping $\boldsymbol{\phi}$ to its unique canonical representative:
    \begin{equation}
    \label{eq:canonical_rep}
    \begin{aligned}
    \boldsymbol{\phi}_{\mathrm{can}} &= \delta(\boldsymbol{\phi}) \cdot \boldsymbol{\phi} \;=\; \bigl(0,\;\; (\phi_{30} - \phi_7) \bmod 7,\;\; 0\bigr) \\
    &\in \mathcal{R}_7 \;:=\; \bigl\{(0, d, 0) : d \in \{0,1,2,3,4,5,6\}\bigr\}.
    \end{aligned}
    \end{equation}
    In particular, within the 210-tuple slice $\{(0, d, \phi_{90}) : d \in \mathbb{Z}_7, \phi_{90} \in \mathbb{Z}_{30}\}$, setting $\phi_7 = 0$ in Eq.~\eqref{eq:crt_shift} gives $\delta = 91\phi_{90} \bmod 210$, showing that all 30 values of $\phi_{90} \in \mathbb{Z}_{30}$ belong to the exact same $\mathbb{Z}_{210}$-orbit as $(0, d, 0)$.
    \item \textbf{Exact Constraint-Matrix Equivalence:} For any $\boldsymbol{\phi} \in \Phi$ and $\delta \in \mathbb{Z}_{210}$, $n_f(s, h; \delta \cdot \boldsymbol{\phi}) = n_f((s + \delta) \bmod 210, h; \boldsymbol{\phi})$ for all $f \in \{7,30,90\}$ and $h \ge 1$. Because $s \mapsto (s + \delta) \bmod 210$ is a bijection on $\mathbb{Z}_{210}$, the multiset of 210 start-day intersection profiles $\{(n_7(s,h), n_{30}(s,h), n_{90}(s,h))\}_{s \in \mathbb{Z}_{210}}$, the maximum exposure $L_{\max}(h; \boldsymbol{\phi}, \boldsymbol{v})$, and the feasible budget polytope $\mathcal{V}(\boldsymbol{\phi})$ are \textbf{identical} across each orbit.
\end{enumerate}
\end{proposition}
\begin{proof}
(1) An element $\delta \in \mathbb{Z}_{210}$ fixes $(\phi_7, \phi_{30}, \phi_{90})$ iff $\delta \equiv 0 \pmod 7$ and $\delta \equiv 0 \pmod{30}$, which by $\operatorname{lcm}(7, 30) = 210$ holds iff $\delta \equiv 0 \pmod{210}$, so $\mathrm{Stab}_{\mathbb{Z}_{210}}(\boldsymbol{\phi}) = \{0\}$ and $|\mathcal{O}(\boldsymbol{\phi})| = |G|/1 = 210$.
(2) Modulo 7, $120\phi_7 + 91\phi_{90} = (17 \times 7 + 1)\phi_7 + (13 \times 7)\phi_{90} \equiv \phi_7 \pmod 7$. Modulo 30, $120\phi_7 + 91\phi_{90} = (4 \times 30)\phi_7 + (3 \times 30 + 1)\phi_{90} \equiv \phi_{90} \pmod{30}$. By the Chinese Remainder Theorem, $\delta(\boldsymbol{\phi})$ in Eq.~\eqref{eq:crt_shift} is the unique residue in $\mathbb{Z}_{210}$ with $(\phi_7 - \delta) \bmod 7 = 0$ and $(\phi_{90} - \delta) \bmod 30 = 0$, while the middle coordinate becomes $(\phi_{30} - \delta) \bmod 7 = (\phi_{30} - \phi_7) \bmod 7$. Since the 7 representatives $\{(0, d, 0) : d \in \mathbb{Z}_7\}$ have distinct values of $(\phi_{30} - \phi_7) \bmod 7$, which is invariant under $\delta \cdot \boldsymbol{\phi}$, they are pairwise non-equivalent and exhaust all 7 orbits.
(3) In Eq.~\eqref{eq:k_bounds}, replacing $\phi_f$ by $(\phi_f - \delta) \bmod s_f = \phi_f - \delta + m_f s_f$ ($m_f \in \mathbb{Z}$) shifts both $k_{\min}$ and $k_{\max}$ by $-m_f$ while replacing $s - \phi_f$ with $(s + \delta) - \phi_f$. Thus $n_f(s, h; \delta \cdot \boldsymbol{\phi}) = n_f(s + \delta, h; \boldsymbol{\phi}) = n_f((s + \delta) \bmod 210, h; \boldsymbol{\phi})$ by Theorem~\ref{thm:periodicity}. Since $\mathbb{Z}_{210} + \delta = \mathbb{Z}_{210}$, maximizing over $s \in \mathbb{Z}_{210}$ is invariant.
\end{proof}

\noindent\textbf{When Is \texorpdfstring{$\mathbb{Z}_{210}$}{Z210} Orbit Equivalence Broken?}
Proposition~\ref{prop:orbits} holds whenever both the privacy constraints and the utility objective depend only on the shift-invariant steady-state schedule. Equivalence is broken in three specific operational regimes:
\begin{enumerate}[leftmargin=*, nosep, topsep=1pt]
    \item \textbf{External Fixed Calendar-Date / Weekday Queries:} If an executive dashboard queries every Sunday ($q = 7k + 6$), the query-date staleness $\tau_f(q; \phi_f) = (q - \phi_f) \bmod s_f$ depends directly on $\phi_{30} \bmod 7$ and $\phi_{90} \bmod 30$ relative to Sunday. For example, $\boldsymbol{\phi} = (0,0,0)$ has Sunday staleness $6$ days on $W_{30}$, whereas $\boldsymbol{\phi} = (0,6,0)$ publishes $W_{30}$ directly on Sundays ($\tau_{30} = 0$).
    \item \textbf{Origin-Anchored Staleness Penalties ($\lambda_{\text{stale}} > 0$):} If the planner objective adds a penalty anchored to fixed calendar dates (or finite startup boundary effects on days $t \in [1, W_f-1]$), shifting by $\delta$ alters the boundary objective.
    \item \textbf{Pinned Phase Constraints (\texttt{fixed\_phi\_7}, \texttt{fixed\_phi\_30}, \texttt{fixed\_phi\_90}):} When business rules pin one or more phase coordinates, the allowed symmetry group shrinks to the subgroup fixing the pinned coordinates; \texttt{JointCalendarPlanner} automatically disables full $\mathbb{Z}_{210}$ orbit reduction when any phase coordinate is pinned.
\end{enumerate}

\subsection{Separable Convex Budget Optimization and \texttt{DECIMAL100} Dual Certificates}
For each fixed canonical phase $\boldsymbol{\phi} \in \mathcal{R}_7$, optimizing the budget vector $\boldsymbol{v} = (d, x_7, x_{30}, x_{90})^\top \in [\boldsymbol{l}, \boldsymbol{u}] \subset \mathbb{R}_{>0}^4$ to minimize the weighted variance proxy is a strictly convex separable program over a finite linear polytope:
\begin{equation}
\label{eq:convex_subproblem}
\min_{\boldsymbol{v} \in [\boldsymbol{l}, \boldsymbol{u}]} \; F(\boldsymbol{v}) := \sum_{j=1}^4 \frac{c_j}{v_j} \qquad \text{s.t.} \qquad A(\boldsymbol{\phi}) \boldsymbol{v} \;\le\; \boldsymbol{b},
\end{equation}
where row 0 of $A(\boldsymbol{\phi}) \boldsymbol{v} \le \boldsymbol{b}$ encodes the asymptotic slope bound $\mathrm{slope}_{\infty}(\boldsymbol{v}) \le \beta_{\mathrm{tail}}$ and rows $1, \dots, 575$ encode $\max_{s \in [0, 209]} (h d + \sum_f n_f(s, h; \phi_f) x_f) \le \bar{\rho}(h)$ for $h \in [1, 575]$. When only a single linear constraint $\sum_{j=1}^4 a_j v_j \le b$ is active in the interior of $[\boldsymbol{l}, \boldsymbol{u}]$, the KKT conditions yield the closed-form square-root allocation $v_j^\star = \frac{b \sqrt{c_j / a_j}}{\sum_{m=1}^4 \sqrt{a_m c_m}}$. Under general multi-horizon constraints where multiple horizons ($h \in \{1, 2, 7, 30, 365\}$ and the tail slope) are simultaneously active, \texttt{JointCalendarPlanner} solves Eq.~\eqref{eq:convex_subproblem} via active-set \texttt{SLSQP}, applies inward ray shrinkage $\boldsymbol{v}_{\text{cert}} = \gamma \hat{\boldsymbol{v}}$ to guarantee strict zero-tolerance feasibility ($\min_{h} \mathrm{slack}(h) \ge 0$), and certifies optimality via 100-digit directed-rounding Lagrangian duality.

\begin{proposition}[Certified Directed-Rounding Lagrangian Dual Lower Bound]
\label{prop:dual}
For any dual vector $\boldsymbol{\lambda} \in \mathbb{R}_{\ge 0}^{576}$ and $\boldsymbol{g} = A(\boldsymbol{\phi})^\top \boldsymbol{\lambda} \ge \mathbf{0}$, weak duality implies that for every feasible $\boldsymbol{v}$ in Eq.~\eqref{eq:convex_subproblem}:
\begin{equation}
\label{eq:dual_lb}
\begin{aligned}
F(\boldsymbol{v}) &\ge q(\boldsymbol{\lambda}) := -\boldsymbol{\lambda}^\top \boldsymbol{b} + \sum_{j=1}^4 \min_{v_j \in [l_j, u_j]} \!\left(\frac{c_j}{v_j} + g_j v_j\right), \\
\hat{v}_j(g_j) &= \mathrm{clip}\!\left(\sqrt{c_j / g_j},\, l_j,\, u_j\right).
\end{aligned}
\end{equation}
Evaluating $q(\boldsymbol{\lambda})$ in IEEE-754 float64 with a heuristic deduction $q_{\text{float64}}(\boldsymbol{\lambda})(1 - 10^{-14})$ is \textbf{unsound}: under large dual multipliers ($\|\boldsymbol{\lambda}\|_1 \sim 10^{15}$), catastrophic cancellation between $\sum_j g_j \hat{v}_j$ and $\boldsymbol{\lambda}^\top \boldsymbol{b}$ causes $q_{\text{float64}}(\boldsymbol{\lambda})(1 - 10^{-14})$ to exceed the true primal optimum by $> 0.19$. Instead, evaluating $g_{j,\mathrm{lo}}$ via \texttt{ROUND\_FLOOR}, $d_{\mathrm{hi}} = \boldsymbol{\lambda}^\top \boldsymbol{b}$ via \texttt{ROUND\_CEILING}, and $q_{\mathrm{dec,lo}}$ in 100-digit decimal arithmetic (\texttt{DECIMAL100\_DIRECTED\_ROUNDING}) followed by downward \texttt{math.nextafter} conversion to float64 yields a mathematically rigorous lower bound $q_{\mathrm{cert}}(\boldsymbol{\lambda}) \le F(\boldsymbol{v}^\star)$.
\end{proposition}
\begin{proof}
Since $\boldsymbol{v} > \mathbf{0}$ and $\boldsymbol{\lambda} \ge \mathbf{0}$, replacing $A(\boldsymbol{\phi})_{i,j}$ by lower bounds $A_{\mathrm{lo},i,j} \le A(\boldsymbol{\phi})_{i,j}$ (including $\lfloor 1/7 \rfloor_{100}$ and $\lfloor 1/30 \rfloor_{100}$ in the tail row) and rounding downward ensures $0 \le g_{j,\mathrm{lo}} \le g_j$. For any candidate $\tilde{v}_j \in [l_j, u_j]$ (and by 1D convexity of $v_j \mapsto c_j/v_j + g_{j,\mathrm{lo}} v_j$, taking the minimum over the two endpoints of the 1-ULP decimal enclosure $[\tilde{v}_{j,\mathrm{lo}}, \tilde{v}_{j,\mathrm{hi}}] \ni \hat{v}_j(g_{j,\mathrm{lo}})$ evaluated with \texttt{ROUND\_FLOOR}), we obtain $\phi_{j,\mathrm{lo}} \le \min_{v_j \in [l_j, u_j]}(c_j/v_j + g_{j,\mathrm{lo}} v_j) \le \min_{v_j \in [l_j, u_j]}(c_j/v_j + g_j v_j)$. Subtracting the upward-rounded $d_{\mathrm{hi}} \ge \boldsymbol{\lambda}^\top \boldsymbol{b}$ under \texttt{ROUND\_FLOOR} guarantees $q_{\mathrm{cert}}(\boldsymbol{\lambda}) \le q(\boldsymbol{\lambda}) \le F(\boldsymbol{v}^\star)$.
\end{proof}

\noindent\textbf{Four-Tier Certification Taxonomy.}
To prevent conflating numerical solver convergence with symbolic proof, \texttt{Joint\allowbreak Calendar\allowbreak Planner} enforces four strict tiers: (1) \texttt{SOLVER\_\allowbreak CONVERGED} (SLSQP converged, uncertified); (2) \texttt{FEASIBLE\_\allowbreak CERTIFIED} (passes 575-day + tail zero-tolerance check after shrinkage); (3) \texttt{NEAR\_\allowbreak GLOBAL\_\allowbreak OPTIMAL\_\allowbreak CERTIFIED} (feasible-certified plus all 7 canonical phase orbits certified with relative primal--dual gap $\eta_{\text{gap}} = (F_{\text{UB}} - LB_{\text{global}})/LB_{\text{global}} \le 10^{-4}$; empirically $\eta_{\text{gap}} \le 7.45 \times 10^{-9}$); and (4) \texttt{GLOBAL\_\allowbreak OPTIMAL\_\allowbreak PROVED} (reserved exclusively for symbolic closed-form proofs; never set by numerical runs).

\begin{table*}[t]
\centering
\caption{Planning benchmark across 4 groups (A: Historical Baseline $d=2.416\times 10^{-4}, x_f=1.530\times 10^{-3}$; B: Budget-Only at $\boldsymbol{\phi}=(0,0,0)$; C: Phase-Only with historical budget; D: Joint Phase + Budget Planning) under both commitments and two workloads (\texttt{Default} $(0.40, 0.25, 0.25, 0.10)$ and \texttt{DistinctHeavy} $(0.10, 0.40, 0.40, 0.10)$). Compares 7-orbit canonical search ($\mathcal{R}_7$) vs.\ 210-tuple enumeration ($\{0\}\times\mathbb{Z}_7\times\mathbb{Z}_{30}$), showing identical certified objectives and \texttt{DECIMAL100} dual bounds with a $19\times\text{--}23\times$ speedup.}
\label{tab:planner_benchmark}
\resizebox{\textwidth}{!}{%
\begin{tabular}{@{}lllccccccccc@{}}
\toprule
\textbf{Curve} & \textbf{Workload} & \textbf{Group} & \textbf{Status / Tier} & \textbf{Phase} $\boldsymbol{\phi}^\star$ & $d$ ($\times 10^{-4}$) & $(x_7, x_{30}, x_{90})$ ($\times 10^{-3}$) & \textbf{Primal UB} $F(\boldsymbol{v}_{\text{cert}})$ & \textbf{\texttt{DECIMAL100} Dual LB} & \textbf{Rel Gap} $\eta_{\text{gap}}$ & \textbf{Min Slack} ($h^\star$) & \textbf{Time (7-Orbit / 210-Slice)} \\ \midrule
\multirow{4}{*}{\texttt{Week1\_v2}} & \multirow{4}{*}{\texttt{Default}} & A\_Baseline & \texttt{FEASIBLE\_CERT} & $(0,0,0)$ & $2.4160$ & $(1.530, 1.530, 1.530)$ & $2047.7860$ & --- & --- & $+1.408 \times 10^{-3}$ ($2$) & $0.020\,\text{s}$ \\
 & & B\_BudgetOnly & \texttt{NEAR\_GLOBAL\_CERT} & $(0,0,0)$ & $4.7113$ & $(0.985, 0.985, 1.290)$ & $\mathbf{1433.9182}$ & $\mathbf{1433.9182}$ & $1.01 \times 10^{-11}$ & $+1.809 \times 10^{-3}$ ($212$) & $0.026\,\text{s}$ \\
 & & C\_PhaseOnly & \texttt{FEASIBLE\_CERT} & $(0,3,0)$ & $2.4160$ & $(1.530, 1.530, 1.530)$ & $2047.7860$ & --- & --- & $+1.853 \times 10^{-3}$ ($32$) & $0.026\,\text{s} \;/\; 0.358\,\text{s}$ \\
 & & D\_JointPlan & \texttt{NEAR\_GLOBAL\_CERT} & $(0,0,0)$ & $4.7113$ & $(0.985, 0.985, 1.290)$ & $\mathbf{1433.9182}$ & $\mathbf{1433.9182}$ & $2.04 \times 10^{-11}$ & $+1.809 \times 10^{-3}$ ($212$) & $\mathbf{0.066\,\text{s}} \;/\; 1.537\,\text{s}$ ($23.2\times$) \\ \midrule
\multirow{3}{*}{\texttt{Week1\_v2}} & \multirow{3}{*}{\texttt{DistHeavy}} & A\_Baseline & \texttt{FEASIBLE\_CERT} & $(0,0,0)$ & $2.4160$ & $(1.530, 1.530, 1.530)$ & $1002.1426$ & --- & --- & $+1.408 \times 10^{-3}$ ($2$) & $0.019\,\text{s}$ \\
 & & B\_BudgetOnly & \texttt{NEAR\_GLOBAL\_CERT} & $(0,0,0)$ & $3.1700$ & $(1.579, 1.442, 1.156)$ & $932.6324$ & $932.6324$ & $1.00 \times 10^{-11}$ & $+8.722 \times 10^{-13}$ ($93$) & $0.029\,\text{s}$ \\
 & & D\_JointPlan & \texttt{NEAR\_GLOBAL\_CERT} & $(0,4,0)$ & $3.1300$ & $(1.611, 1.459, 1.232)$ & $\mathbf{922.9791}$ & $\mathbf{922.9791}$ & $7.45 \times 10^{-9}$ & $+9.335 \times 10^{-13}$ ($98$) & $\mathbf{0.068\,\text{s}} \;/\; 1.509\,\text{s}$ ($22.4\times$) \\ \midrule
\multirow{3}{*}{\texttt{Original}} & \multirow{3}{*}{\texttt{Default}} & A\_Baseline & \texttt{CALIB\_FALLBACK} & $(0,0,0)$ & $2.1237$ & $(1.345, 1.345, 1.345)$ & $2329.5959$ & --- & --- & $+1.522 \times 10^{-13}$ ($2$) & $0.022\,\text{s}$ \\
 & & B\_BudgetOnly & \texttt{NEAR\_GLOBAL\_CERT} & $(0,0,0)$ & $4.7000$ & $(0.983, 0.983, 1.287)$ & $\mathbf{1437.6283}$ & $\mathbf{1437.6283}$ & $1.02 \times 10^{-11}$ & $+1.298 \times 10^{-3}$ ($212$) & $0.022\,\text{s}$ \\
 & & D\_JointPlan & \texttt{NEAR\_GLOBAL\_CERT} & $(0,3,0)$ & $4.7000$ & $(0.983, 0.983, 1.287)$ & $\mathbf{1437.6283}$ & $\mathbf{1437.6283}$ & $2.01 \times 10^{-11}$ & $+1.579 \times 10^{-3}$ ($122$) & $\mathbf{0.050\,\text{s}} \;/\; 0.952\,\text{s}$ ($19.1\times$) \\ \midrule
\multirow{3}{*}{\texttt{Original}} & \multirow{3}{*}{\texttt{DistHeavy}} & A\_Baseline & \texttt{CALIB\_FALLBACK} & $(0,0,0)$ & $2.1237$ & $(1.345, 1.345, 1.345)$ & $1140.0543$ & --- & --- & $+1.522 \times 10^{-13}$ ($2$) & $0.067\,\text{s}$ \\
 & & B\_BudgetOnly & \texttt{NEAR\_GLOBAL\_CERT} & $(0,0,0)$ & $3.2100$ & $(1.574, 1.411, 1.093)$ & $940.8881$ & $940.8881$ & $1.00 \times 10^{-11}$ & $+1.522 \times 10^{-13}$ ($2$) & $0.067\,\text{s}$ \\
 & & D\_JointPlan & \texttt{NEAR\_GLOBAL\_CERT} & $(0,4,0)$ & $3.1500$ & $(1.605, 1.435, 1.175)$ & $\mathbf{930.3260}$ & $\mathbf{930.3260}$ & $1.02 \times 10^{-11}$ & $+8.807 \times 10^{-13}$ ($95$) & $\mathbf{0.081\,\text{s}} \;/\; 1.701\,\text{s}$ ($21.0\times$) \\ \bottomrule
\end{tabular}%
}
\end{table*}

\section{Experimental Evaluation}
\label{sec:experiments}

\subsection{Experimental Setup and Statistical Methodology}
\noindent\textbf{420-Day Cross-Period Simulation \& Collision-Free 64-Bit CRN Seeding.}
To evaluate SDG-MH across two full 210-day super-periods ($T = 420 = 2P$ days, exceeding $H_0 = 365$), we generate 10 independent enterprise telemetry streams (data seeds $s_d \in \{1042, 2042, \dots, 10042\}$, $N=180$ users, $M=9$ cells, heavy-tailed user activity with weekly seasonality). For each data seed and each of the 16 configurations ($2\text{ curves} \times 2\text{ workloads} \times 4\text{ groups}$), we execute 10 Monte Carlo (MC) trials ($1{,}600$ total 420-day runs). To prevent the 32-bit modular seed collision identified in earlier linear-combination seeding (where $(3042,2,410,1)$ and $(5042,4,218,2)$ collided at $3{,}129{,}808{,}095$), each channel draw uses an injective 64-bit pack $\mathrm{pack64}(s_d, r, t, c) = (s_d \ll 36) \mid (r \ll 24) \mid (t \ll 8) \mid c$ passed into \texttt{np.random.SeedSequence}, verified collision-free across all $10 \times 10 \times 420 \times 5 = 210{,}000$ channel draws. Crucially, $(s_d, r, t, c)$ depends \emph{only} on data seed $s_d$, trial $r$, day $t$, and channel $c$ (never on group or curve), so on common release days Groups A/B/C/D share identical standard normal draws $Z_{s_d,r,t,c} \sim \mathcal{N}(0, I)$, isolating the exact effect of budget and phase optimization.

\noindent\textbf{Two-Stage Hierarchical Pairing \& Post-Processing Consistency.}
Because 10 MC draws on the same telemetry stream share underlying user trajectories and overlapping windows, treating $3 \times 10 = 30$ or $10 \times 10 = 100$ trials as independent i.i.d.\ samples understates uncertainty. We first average across the 10 MC trials within each data seed $s_d$ to form $S=10$ strictly independent seed-level means $\bar{Y}_{s_d}$, and compute two-sided 95\% Student-$t$ confidence intervals ($df=9$, $t_{0.975,9}=2.2622$), Wilcoxon signed-rank tests, sign tests, and Holm--Bonferroni family-wise adjusted $p$-values ($m=12$). Across all mechanisms (SDG-MH and continual baselines), post-processing is unified: $\hat{R}_t = \max(0, \tilde{R}_t)$ and $\hat{F}_t = \min(\hat{R}_t, \max(0, \tilde{F}_t))$.

\subsection{Joint Planning vs.\ Budget-Only: Same-Date Accuracy, Staleness, and Truncation Bias}
Table~\ref{tab:paired_bd} compares Group D (joint planning) and Group B (budget-only) across all $S=10$ data seeds. Three findings stand out:

\begin{enumerate}[leftmargin=*, nosep, topsep=2pt]
    \item \textbf{Certified Variance Reduction:} Under \texttt{Default} weights ($w_{\text{daily}}=0.40$), the asymptotic tail constraint dominates so $B$ and $D$ achieve the same objective ($1433.92$ and $1437.63$), whereas under \texttt{DistinctHeavy} weights ($w_7=w_{30}=0.40$), phase staggering $\boldsymbol{\phi}^\star=(0,4,0)$ eliminates multi-window collision peaks at $h \in \{2, 93\}$, reducing certified objective $F(\boldsymbol{v}_{\text{cert}})$ from $932.63 \to 922.98$ (\texttt{Week1\_v2}) and $940.89 \to 930.33$ (\texttt{Original\_Planner}), and cutting raw Gaussian noise variance $\sigma_{\text{Distinct30}}^2$ by $-20.05\%$ ($6{,}334.7 \to 5{,}064.4$).
    \item \textbf{Why Lower Noise Variance Can Increase Distinct MAPE (Negative Truncation Bias):} Despite a $-20.05\%$ reduction in Distinct30 Gaussian noise variance on \texttt{Original\_Planner / DistinctHeavy}, release-day Distinct30 MAPE increases by $+4.03\%$ ($35.61\% \to 39.65\%$, $p_{\text{Holm}}=0.054$). Why? Window-level distinct counting clips multi-cell user indicators to $\Delta_2 = \sqrt{11}$, inducing a deterministic \emph{negative clipping bias} $\mu_{\text{clip}} \approx -39.95\%$ relative to true unclipped distinct users, while post-processing applies $\max(0, \cdot)$ truncation. For a negatively biased estimator $Y = (\mu + b) + \sigma Z$ with $b < 0$ and $\mu + b > 0$, positive zero-truncation $\max(0, Y)$ clips away negative Gaussian realizations and pushes $\mathbb{E}[\max(0, Y)]$ \emph{upward} toward the unclipped truth $\mu$:
    \[
    \mathbb{E}[\max(0, Y)] \;=\; (\mu + b)\Phi\!\left(\frac{\mu+b}{\sigma}\right) + \sigma \phi\!\left(\frac{\mu+b}{\sigma}\right).
    \]
    Consequently, when Group D reduces $\sigma$, fewer negative tails are truncated at zero, reducing the upward zero-truncation bounce and moving noisy releases closer to the \emph{clipped} target (where RMSE improves), which slightly widens the gap to the \emph{unclipped} ground truth!
    \item \textbf{Same-Query-Date Latest-Available Accuracy vs.\ Staleness:} Because Group B ($\boldsymbol{\phi}=(0,0,0)$) and Group D ($\boldsymbol{\phi}=(0,4,0)$) publish $W_{30}$ on different weekdays, comparing only their own release days conflates date effects. Evaluating all $420$ daily queries (and Sunday-pinned queries) using the latest available release shows that Group D maintains identical $3.00$-day mean staleness across all 420 days while shifting Sunday-specific staleness from $0.00$ days to $5.00\text{--}6.00$ days when unpinned, confirming our symmetry-breaking analysis in Section~\ref{sec:planner}: if an application requires zero staleness on a specific weekday, that phase coordinate must be pinned via \texttt{fixed\_phi\_30}.
\end{enumerate}

\begin{table*}[t]
\centering
\caption{Paired comparison ($D - B$: \texttt{JointPlanning} minus \texttt{BudgetOnly}) across $S=10$ independent 420-day data seeds ($\times 10$ CRN trials; 100 paired runs per cell). Reports both own-release-day accuracy and same-query-date latest-available accuracy across all 420 days, with 95\% $t$-CIs ($df=9$) and Holm--Bonferroni adjusted $p$-values ($m=12$).}
\label{tab:paired_bd}
\resizebox{\textwidth}{!}{%
\begin{tabular}{@{}llccccccccc@{}}
\toprule
\textbf{Commitment} & \textbf{Workload} & \textbf{Metric} & \textbf{Channel} & \textbf{Mean $B$} & \textbf{Mean $D$} & \textbf{Paired $\Delta(D-B)$} & \textbf{95\% CI ($df=9$)} & \textbf{Raw $p$} & \textbf{Holm $p$} & \textbf{Wins ($D<B$)} \\ \midrule
\multirow{6}{*}{\texttt{Week1\_v2}} & \multirow{3}{*}{\texttt{Default}} & Own-Release MAPE (\%) & Requests (Daily) & $10.03$ & $10.22$ & $+0.193$ & $[+0.191, +0.194]$ & $<10^{-4}$ & $<10^{-4}$ & $0/10$ \\
 & & Own-Release MAPE (\%) & Distinct 7d & $21.22$ & $19.99$ & $\mathbf{-1.235}$ & $[-1.249, -1.220]$ & $<10^{-4}$ & $<10^{-4}$ & $\mathbf{10/10}$ \\
 & & Own-Release MAPE (\%) & Distinct 30d & $42.06$ & $43.92$ & $+1.863$ & $[-1.608, +5.334]$ & $0.256$ & $0.767$ & $3/10$ \\ \cmidrule(l){2-11}
 & \multirow{3}{*}{\texttt{DistHeavy}} & Own-Release MAPE (\%) & Requests (Daily) & $14.70$ & $14.84$ & $+0.132$ & $[+0.131, +0.134]$ & $<10^{-4}$ & $<10^{-4}$ & $0/10$ \\
 & & Own-Release MAPE (\%) & Distinct 7d & $19.44$ & $17.19$ & $\mathbf{-2.248}$ & $[-2.269, -2.226]$ & $<10^{-4}$ & $<10^{-4}$ & $\mathbf{10/10}$ \\
 & & Own-Release MAPE (\%) & Distinct 30d & $35.42$ & $39.87$ & $+4.448$ & $[+1.274, +7.621]$ & $0.011$ & $0.046$ & $1/10$ \\ \midrule
\multirow{6}{*}{\texttt{Original}} & \multirow{3}{*}{\texttt{Default}} & Own-Release MAPE (\%) & Requests (Daily) & $10.11$ & $10.22$ & $+0.112$ & $[+0.111, +0.113]$ & $<10^{-4}$ & $<10^{-4}$ & $0/10$ \\
 & & Own-Release MAPE (\%) & Distinct 7d & $20.90$ & $20.06$ & $\mathbf{-0.838}$ & $[-0.850, -0.827]$ & $<10^{-4}$ & $<10^{-4}$ & $\mathbf{10/10}$ \\
 & & Own-Release MAPE (\%) & Distinct 30d & $41.61$ & $43.82$ & $+2.203$ & $[-1.264, +5.670]$ & $0.184$ & $0.736$ & $3/10$ \\ \cmidrule(l){2-11}
 & \multirow{3}{*}{\texttt{DistHeavy}} & Own-Release MAPE (\%) & Requests (Daily) & $14.73$ & $14.84$ & $+0.111$ & $[+0.109, +0.112]$ & $<10^{-4}$ & $<10^{-4}$ & $0/10$ \\
 & & Own-Release MAPE (\%) & Distinct 7d & $19.47$ & $17.27$ & $\mathbf{-2.205}$ & $[-2.226, -2.185]$ & $<10^{-4}$ & $<10^{-4}$ & $\mathbf{10/10}$ \\
 & & Own-Release MAPE (\%) & Distinct 30d & $35.61$ & $39.65$ & $+4.034$ & $[+0.880, +7.189]$ & $0.018$ & $0.054$ & $1/10$ \\ \midrule
\multicolumn{2}{@{}l}{\texttt{Week1\_v2 / Default}} & Same-Date 420d MAPE & Distinct 30d (Latest) & $41.32$ & $43.36$ & $+2.043$ & $[-1.160, +5.247]$ & $0.183$ & --- & Staleness: $3.00\text{d}$ vs $3.01\text{d}$ \\
\multicolumn{2}{@{}l}{\texttt{Week1\_v2 / DistHeavy}} & Same-Date 420d MAPE & Distinct 30d (Latest) & $35.33$ & $39.06$ & $+3.733$ & $[+0.785, +6.680]$ & $0.019$ & --- & Staleness: $3.00\text{d}$ vs $3.00\text{d}$ \\ \bottomrule
\end{tabular}%
}
\end{table*}

\subsection{Comparison with Continual Observation Baselines}
Table~\ref{tab:baselines} compares SDG-MH against three continual observation baselines under identical per-day user contribution clipping ($C_{\text{req}}=50, C_{\text{fail}}=10$) and identical post-processing ($\hat{R}_t = \max(0, \tilde{R}_t)$, $\hat{F}_t = \min(\hat{R}_t, \max(0, \tilde{F}_t))$):
\begin{itemize}[leftmargin=*, nosep, topsep=2pt]
    \item \textbf{Binary Tree (\texttt{Tree-2l}, Honaker 2-Level~\cite{honaker2015efficient}) \& \texttt{Tree-U} (Uniform Level~\cite{dwork2010continual,chan2011private}):} Configured over horizon $T=420$ ($L=10$ tree levels) calibrated to the finite 365-day additive budget $\bar{\rho}_{\text{add}}(365) = 0.088267$.
    \item \textbf{$w$-Event Budget Division (\texttt{W-BD-30}~\cite{kellaris2014wevent}):} Uniformly divides the 30-day additive budget $\bar{\rho}_{\text{add}}(30) = 0.011233$ across $w=30$ days ($d = 0.00037444/\text{day}$).
\end{itemize}
Table~\ref{tab:baselines} highlights two critical architectural trade-offs: (1) on short-to-medium additive windows ($h \in \{1, 7, 30\}$), SDG-MH achieves $2.4\times\text{--}8.1\times$ lower MAPE than \texttt{Tree-2l} and \texttt{Tree-U} because binary trees spread their budget across $L=10$ tree levels and violate SDG-MH's short-horizon commitments $\bar{\rho}(1)$ and $\bar{\rho}(7)$ by $5.0\times$ and $3.5\times$; (2) on long additive intervals ($h=365$), \texttt{Tree-2l} achieves lower additive error ($2.09\%$ vs.\ $3.03\%$) because dyadic decomposition reconstructs 365-day sums from at most $7$ nodes ($O(\log^3 h)$ variance) versus SDG-MH's 365 daily independent Gaussians ($O(h)$ variance); and (3) neither \texttt{Tree-2l}, \texttt{Tree-U}, nor \texttt{W-BD-30} can answer non-additive rolling 7/30/90-day distinct-user queries or reserve budget for them.

\begin{table*}[t]
\centering
\caption{Comparison of SDG-MH (\texttt{Week1\_v2 / Default}, Group B and Group D) against continual observation baselines over $T=420$ days under unified post-processing ($\hat{R}_t=\max(0,\tilde{R}_t)$, $\hat{F}_t=\min(\hat{R}_t,\max(0,\tilde{F}_t))$). Strictly separates infinite-horizon multi-knot compliance from finite-horizon $T=420$ tree calibration.}
\label{tab:baselines}
\resizebox{\textwidth}{!}{%
\begin{tabular}{@{}lcccccccccc@{}}
\toprule
\textbf{Mechanism} & \textbf{Horizon Scope} & \textbf{Supports Distinct?} & \textbf{Multi-Knot $\bar{\rho}(h)$ Safe?} & \textbf{Req $h=1$ MAPE} & \textbf{Req $h=7$ MAPE} & \textbf{Req $h=30$ MAPE} & \textbf{Req $h=90$ MAPE} & \textbf{Req $h=365$ MAPE} & \textbf{Fail $h=1$ MAPE} & \textbf{Distinct7 / 30 / 90 MAPE} \\ \midrule
\textbf{SDG-MH (Group B)} & Infinite ($t \in [1, \infty)$) & \textbf{Yes} ($W_{7,30,90}$) & \textbf{Yes} ($\forall h \ge 1$) & $\mathbf{10.03\%}$ & $\mathbf{3.82\%}$ & $1.91\%$ & $1.10\%$ & $3.00\%$ & $\mathbf{51.82\%}$ & $21.22\% \;/\; 42.06\% \;/\; 53.78\%$ \\
\textbf{SDG-MH (Group D)} & Infinite ($t \in [1, \infty)$) & \textbf{Yes} ($W_{7,30,90}$) & \textbf{Yes} ($\forall h \ge 1$) & $10.22\%$ & $3.89\%$ & $1.94\%$ & $1.12\%$ & $3.03\%$ & $52.53\%$ & $\mathbf{19.99\%} \;/\; 43.92\% \;/\; 56.11\%$ \\ \midrule
\textbf{W-BD-30} ($w$-Event~\cite{kellaris2014wevent}) & Infinite ($w=30$ only) & No (Additive only) & Violates $h=365$ & $9.94\%$ & $3.78\%$ & $\mathbf{1.89\%}$ & $1.09\%$ & $2.98\%$ & $51.49\%$ & N/A (Unbudgeted) \\
\textbf{Tree-2l} (Honaker~\cite{honaker2015efficient}) & Finite ($T=420, L=10$) & No (Additive only) & Violates $h \in \{1,7,30\}$ & $50.09\%$ & $20.66\%$ & $7.19\%$ & $2.62\%$ & $\mathbf{2.09\%}$ & $98.06\%$ & N/A (Unbudgeted) \\
\textbf{Tree-U} (Chan et al.~\cite{chan2011private}) & Finite ($T=420, L=10$) & No (Additive only) & Violates $h \in \{1,7,30\}$ & $79.96\%$ & $31.06\%$ & $10.17\%$ & $3.79\%$ & $3.26\%$ & $111.78\%$ & N/A (Unbudgeted) \\ \bottomrule
\end{tabular}%
}
\end{table*}

\section{Limitations, Reviewer Risk Assessment, and Future Benchmark Protocols}
\label{sec:limitations}

To maintain complete scientific transparency, we document three primary limitations of the current study alongside concrete, reproducible experimental protocols designed for follow-up validation:

\noindent\textbf{1. Real-World Telemetry Evaluation Protocol (WildChat-1M).}
Because repository governance (\texttt{DATA\_APPROVAL.md} absence) prohibits accessing internal production logs, Section~\ref{sec:experiments} uses synthetic heavy-tailed enterprise traces. While synthetic traces enable exact ground-truth sensitivity and 10-seed $\times$ 10-trial CRN pairing, human reviewers may question whether real LLM burstiness and non-stationary user churn alter distinct-user clipping bias. \emph{Minimal Reproducible Real-Data Protocol:} Map the public \textbf{WildChat-1M} dataset~\cite{zhao2024wildchat} (timestamped real user--LLM conversations with hashed \texttt{user\_id} and model tags) onto our 9-cell schema over a 420-day window: (a) bucket models into 3 provider tiers and classify prompt categories into 3 task buckets; (b) replay the exact SDG-MH Group A/B/C/D schedules under \texttt{WEEK1\_V2\_CURVE} across 10 CRN noise seeds; and (c) measure empirical user-window active-day distributions $|D_{u,B_{f,t}}|$, distinct clipping bias $\mu_{\text{clip}}$ across $C_{\text{dist}} \in \{\sqrt{5}, \sqrt{11}, \sqrt{18}\}$, and same-date MAPE/RMSE.

\noindent\textbf{2. Multi-Horizon Banded / Toeplitz Matrix Factorization Baseline Protocol.}
For additive queries over $h=365$, daily independent Gaussians incur $O(h)$ variance whereas banded/Toeplitz matrix factorizations~\cite{choquette2023banded,henzinger2024group,mcmahan2025blt} achieve $O(\log^3 h)$ variance. \emph{Minimal Reproducible Banded-MF Protocol:} Construct a $b$-banded lower-triangular strategy matrix $C_b \in \mathbb{R}^{T \times T}$ ($b \in \{7, 30, 90\}$) for the daily additive channel: (a) under $b$-banded column structure, a user's contribution on day $t$ affects at most $b$ rows of $C_b X$, so the $h$-day user-time $\ell_2$-sensitivity of $C_b$ over $I_{s,h}$ is bounded by $\Delta_{2,\text{band}}(h) \le \min(h + b - 1, T) \cdot \max_{t} \|C_{b,*,t}\|_2$; (b) calibrate the noise scale of $C_b$ to satisfy the multi-knot additive curve $\bar{\rho}_{\text{add}}(h) = \bar{\rho}(h) - \sum_f n_f(s,h;\phi_f^\star)x_f^\star$ across all $h \in [1, 575]$ simultaneously; and (c) compare additive RMSE at $h \in \{1, 7, 30, 90, 365\}$ against SDG-MH daily sealing while keeping the periodic distinct-user channel identical.

\noindent\textbf{3. Clipping-Bias Mitigation for Distinct Counts.}
As analyzed in Section~\ref{sec:experiments}, fixed $\ell_2$ clipping ($\Delta_2 = \sqrt{11}$) on multi-cell distinct presence vectors introduces a $\approx -40\%$ negative clipping bias that dominates Gaussian variance on 30-day and 90-day windows. Future work should integrate private quantile-calibrated clipping thresholds or debiasing estimators within the multi-horizon zCDP ledger.

\section{Conclusion}
\label{sec:conclusion}
We presented SDG-MH, a perpetual release architecture and certified joint calendar-and-budget planner for Multi-Horizon User-Time zCDP. By combining exact window-intersection accounting, a non-concave 575-day affine-tail reduction, an exact CRT reduction from $1{,}470$ phase combinations (and $210$ $\phi_7=0$ tuples) to 7 canonical phase orbits, and 100-digit directed-rounding Lagrangian dual certificates, SDG-MH bridges formal continual differential privacy and long-running enterprise LLM telemetry.

\appendix

\section{Supplementary Verification Details}
\label{app:proofs}

\subsection{Symbolic and Numerical Verification of the 7 Canonical Orbits}
For $\Phi = \mathbb{Z}_7 \times \mathbb{Z}_7 \times \mathbb{Z}_{30}$ ($1{,}470$ tuples), let $\mathcal{M}(\boldsymbol{\phi})$ denote the sorted multiset of the 210 start-day intersection profiles across $h \in [1, 575]$:
\[
\mathcal{M}(\boldsymbol{\phi}) \;:=\; \Bigl\{\bigl(n_7(s,h;\phi_7),\, n_{30}(s,h;\phi_{30}),\, n_{90}(s,h;\phi_{90})\bigr)_{h=1}^{575}\Bigr\}_{s=0}^{209}.
\]
Exhaustive evaluation in \texttt{test\_joint\_planner.py} confirms that:
\begin{enumerate}[leftmargin=*, nosep, topsep=1pt]
    \item Across all $1{,}470$ tuples in $\Phi$, $\{\mathcal{M}(\boldsymbol{\phi}) : \boldsymbol{\phi} \in \Phi\}$ has cardinality \textbf{exactly 7}, with each equivalence class containing exactly $210$ tuples in $\Phi$ and exactly $30$ tuples in the $\phi_7 = 0$ slice $\{(0, \phi_{30}, \phi_{90})\}$.
    \item For every $\boldsymbol{\phi} \in \Phi$, the CRT shift $\delta(\boldsymbol{\phi}) = (120\phi_7 + 91\phi_{90}) \bmod 210$ maps $\boldsymbol{\phi}$ to $\boldsymbol{\phi}_{\mathrm{can}} = (0, (\phi_{30} - \phi_7) \bmod 7, 0)$ with $\mathcal{M}(\boldsymbol{\phi}) = \mathcal{M}(\boldsymbol{\phi}_{\mathrm{can}})$ identically across all $h \in [1, 575]$.
    \item Table~\ref{tab:orbits_summary} lists all 7 canonical orbit representatives $(0, d, 0)$ for $d \in \mathbb{Z}_7$ and their certified primal and \texttt{DECIMAL100} dual bounds on \texttt{Week1\_v2 / DistinctHeavy}.
\end{enumerate}

\begin{table}[h]
\centering
\caption{Exact values across all 7 canonical phase orbits $\mathcal{R}_7 = \{(0,d,0) : d \in \mathbb{Z}_7\}$ for \texttt{Week1\_v2 / DistinctHeavy} (each representing 210 tuples in $\Phi$ and 30 tuples in $\{0\}\times\mathbb{Z}_7\times\mathbb{Z}_{30}$).}
\label{tab:orbits_summary}
\resizebox{\columnwidth}{!}{%
\begin{tabular}{@{}ccccccc@{}}
\toprule
\textbf{Orbit} $d$ & \textbf{Rep} $\boldsymbol{\phi}_{\text{can}}$ & \textbf{Orbit Size} ($\Phi$ / Slice) & \textbf{Group C Min Slack} ($h^\star$) & \textbf{Group D Primal UB} & \textbf{\texttt{DECIMAL100} Dual LB} & \textbf{Tightest} $h^\star$ \\ \midrule
$d=0$ & $(0,0,0)$ & $210 \;/\; 30$ & $+1.4082 \times 10^{-3}$ ($2$) & $932.6324$ & $932.6324$ & $93$ \\
$d=1$ & $(0,1,0)$ & $210 \;/\; 30$ & $+1.4082 \times 10^{-3}$ ($2$) & $932.6324$ & $932.6324$ & $93$ \\
$d=2$ & $(0,2,0)$ & $210 \;/\; 30$ & $+1.4082 \times 10^{-3}$ ($2$) & $932.6324$ & $932.6324$ & $93$ \\
$d=3$ & $(0,3,0)$ & $210 \;/\; 30$ & $+1.8526 \times 10^{-3}$ ($32$) & $927.5552$ & $927.5552$ & $94$ \\
$d=4$ & $(0,4,0)$ & $210 \;/\; 30$ & $+1.8526 \times 10^{-3}$ ($32$) & $\mathbf{922.9791}$ & $\mathbf{922.9791}$ & $98$ \\
$d=5$ & $(0,5,0)$ & $210 \;/\; 30$ & $+1.8526 \times 10^{-3}$ ($32$) & $\mathbf{922.9791}$ & $\mathbf{922.9791}$ & $98$ \\
$d=6$ & $(0,6,0)$ & $210 \;/\; 30$ & $+1.8526 \times 10^{-3}$ ($32$) & $927.5552$ & $927.5552$ & $94$ \\ \bottomrule
\end{tabular}%
}
\end{table}

\bibliographystyle{ACM-Reference-Format}
\bibliography{references}

\end{document}
